\documentclass[12pt,a4paper]{article}
\usepackage[utf8]{inputenc}
\usepackage[croatian]{babel}
\usepackage[T1]{fontenc}
\usepackage{graphicx}
\usepackage{geometry}
\usepackage{setspace}
\usepackage{titlesec}
\usepackage{fancyhdr}
\usepackage{caption}

% Postavke margina
\geometry{left=3cm,right=2.5cm,top=2.5cm,bottom=2.5cm}

% Prored 1.5
\onehalfspacing

% Stilovi naslova
\titleformat{\section}{\normalfont\Large\bfseries\uppercase}{}{0pt}{}
\titleformat{\subsection}{\normalfont\large\bfseries}{}{0pt}{}

\begin{document}

% --- NASLOVNICA ---
\begin{titlepage}
    \begin{center}
        {\large SVEUČILIŠTE U ZAGREBU} \\
        {\large FAKULTET STROJARSTVA I BRODOGRADNJE} \\
        \vspace{6cm}
        {\large KOLEGIJ: Humanoidna robotika} \\
        \vspace{2cm}
        {\Huge \textbf{SEMINAR}} \\
        \vspace{1cm}
        {\Large \textbf{Detekcija voća i procjena poze primjenom dubokog učenja i 3D vizije}} \\
        \vfill
        {\large Krešimir Hartl} \\
        \vspace{2cm}
        {\large Zagreb, 2026.}
    \end{center}
\end{titlepage}

% --- SADRŽAJ I POPISI ---
\tableofcontents
\newpage
\listoffigures
\newpage
\listoftables
\newpage

% --- SADRŽAJ (Ovdje AI ubacuje tekst) ---

\section{UVOD}
U ovom seminaru obrađuje se implementacija sustava za detekciju i procjenu 3D poze objekata (voća) u prostoru, što je ključan korak u razvoju sustava za robotsku manipulaciju i vizualnu percepciju humanoidnih robota. Rad se temelji na izradi vlastitog skupa podataka, obučavanju modela dubokog učenja za segmentaciju instanci te obradi trodimenzionalnih oblaka točaka (engl. point cloud) dobivenih pomoću RGB-D senzora. Cilj je bio osposobiti sustav za prepoznavanje različitih vrsta voća te točno određivanje njihovog težišta (centroida) u složenim scenama s višestrukim objektima. Uspješno rješavanje ovog problema omogućuje robotu interakciju s okolinom, kao što je prepoznavanje i hvatanje željenog objekta iz posude.

\section{PRIKUPLJANJE I PRIPREMA SKUPA PODATAKA}

\subsection{Postupak prikupljanja slika}
Prvi korak u razvoju algoritma dubokog učenja bila je izrada prilagođenog skupa podataka koji reprezentira stvarne uvjete u okolini robota. Korišteno je 6 klasa voća (npr. crvena jabuka, zelena jabuka, limun, banana, avokado, orah, kruška - u obliku 3D printanih modela). Za svaku klasu prikupljeno je između 100 i 120 fotografija pojedinačnih plodova, pri čemu se vodilo računa o varijaciji kuta snimanja, mjerila, osvjetljenja te pojavi djelomičnih okluzija.

Nadalje, snimljeno je i 80 do 120 slika složenih scena u kojima se više komada različitog voća nalazi u posudama ili na stolu. Na ovim slikama voće se preklapa, čime se simuliraju stvarni uvjeti s kojima se robot susreće pri zadatku sortiranja ili odabira.

\subsection{Označavanje i integracija podataka}
Nakon prikupljanja slika, izvedeno je označavanje (anotacija) na razini piksela, odnosno primijenjena je metoda segmentacije instanci (engl. instance segmentation). Prije početka označavanja, kreirana je zajednička \texttt{.csv} datoteka kojom je standardizirano nazivlje klasa. Svi sudionici u projektu su nezavisno označavali svoj podskup slika, bez primjene prethodne augmentacije. 

Nakon završene anotacije, svi su pojedinačni skupovi podataka integrirani u jedan centralni skup unutar platforme Roboflow. Tek je na spojenom i usklađenom skupu podataka izvedena augmentacija slika, čime je dodatno povećana robusnost konačnog dataseta namijenjenog treniranju konvolucijskih neuronskih mreža.

\begin{figure}[h!]
    \centering
    \includegraphics[width=0.45\textwidth]{slike/dataset_jabuka.png}
    \hfill
    \includegraphics[width=0.45\textwidth]{slike/dataset_kutija.png}
    \caption{Primjer prikupljenih slika: izdvojeni plod (lijevo) i složena scena s preklapanjem voća u kutiji (desno).}
    \label{fig:okluzije}
\end{figure}

\begin{table}[h!]
    \centering
    \caption{Parametri prikupljenog skupa podataka}
    \vspace{0.3cm}
    \begin{tabular}{|l|c|}
    \hline
    \textbf{Parametar} & \textbf{Vrijednost} \\
    \hline
    Broj klasa voća & 6 \\
    Slika po klasi (pojedinačne) & 100 - 120 \\
    Slika složenih scena & 80 - 120 \\
    Tip anotacije & Segmentacija instanci \\
    \hline
    \end{tabular}
    \label{tab:parametri}
\end{table}

\section{TRENING MODELA I DETEKCIJA}
Za zadatak detekcije i segmentacije voća korišten je model iz YOLO arhitekture (YOLO varijanta prilagođena za segmentaciju). Spojeni i augmentirani skup podataka s Roboflow platforme preuzet je i korišten za treniranje. Model je evaluiran i optimiziran kako bi s visokom pouzdanošću detektirao maske pojedinačnih plodova voća, čak i kada su oni djelomično prekriveni drugim objektima. Dobiveni model predstavlja osnovu za daljnje mapiranje slikovnih informacija u trodimenzionalni prostor.

\subsection{Rezultati na validacijskom skupu}

Kako bi se detaljno prikazala kvaliteta naučenog modela, u nastavku su izložene sve metrike i vizualizacije dobivene na validacijskom skupu.

\textbf{Usporedba stvarnih i predviđenih maski:}
Na slikama ispod prikazani su odabrani primjeri iz validacijskog skupa. Lijevo su prikazane stvarne (ručno označene) maske, dok su desno prikazane predikcije modela na tim istim slikama.
\begin{figure}[ht!]
    \centering
    \includegraphics[width=0.45\textwidth]{slike/val_batch0_labels.jpg}
    \hfill
    \includegraphics[width=0.45\textwidth]{slike/val_batch0_pred.jpg}
    \caption{Validacijski Batch 0: Stvarne oznake (lijevo) i predikcije modela (desno).}
\end{figure}

\begin{figure}[ht!]
    \centering
    \includegraphics[width=0.45\textwidth]{slike/val_batch1_labels.jpg}
    \hfill
    \includegraphics[width=0.45\textwidth]{slike/val_batch1_pred.jpg}
    \caption{Validacijski Batch 1: Stvarne oznake (lijevo) i predikcije modela (desno).}
\end{figure}

\begin{figure}[ht!]
    \centering
    \includegraphics[width=0.45\textwidth]{slike/val_batch2_labels.jpg}
    \hfill
    \includegraphics[width=0.45\textwidth]{slike/val_batch2_pred.jpg}
    \caption{Validacijski Batch 2: Stvarne oznake (lijevo) i predikcije modela (desno).}
\end{figure}

\textbf{Matrice zabune (Confusion Matrix):}
Matrica zabune daje uvid u točnost klasifikacije pojedinih klasa. S obzirom na to da su izračunate apsolutna i normalizirana matrica, iz njih je vidljivo u kojoj mjeri model miješa klase voća međusobno, kao i s pozadinom.
\begin{figure}[ht!]
    \centering
    \includegraphics[width=0.45\textwidth]{slike/confusion_matrix.png}
    \hfill
    \includegraphics[width=0.45\textwidth]{slike/confusion_matrix_normalized.png}
    \caption{Matrica zabune u apsolutnim (lijevo) i normaliziranim (desno) vrijednostima.}
\end{figure}

\textbf{Metrike graničnih kutija (Bounding Box):}
Sljedeće krivulje opisuju performanse detekcije objekata pomoću graničnih kutija. Prikazane su krivulje Preciznosti (P), Odziva (R), F1 mjere te odnosa Preciznosti i Odziva (PR krivulja).
\begin{figure}[ht!]
    \centering
    \includegraphics[width=0.45\textwidth]{slike/BoxP_curve.png}
    \hfill
    \includegraphics[width=0.45\textwidth]{slike/BoxR_curve.png}
    
    \vspace{0.5cm}
    \includegraphics[width=0.45\textwidth]{slike/BoxF1_curve.png}
    \hfill
    \includegraphics[width=0.45\textwidth]{slike/BoxPR_curve.png}
    \caption{Box metrike: Preciznost (gore lijevo), Odziv (gore desno), F1 (dolje lijevo) i PR krivulja (dolje desno).}
\end{figure}

\textbf{Metrike segmentacijskih maski (Mask):}
Slično kao i za granične kutije, segmentacija na razini piksela također ima svoje mjere preciznosti i odziva koje su prikazane na donjim krivuljama.
\begin{figure}[ht!]
    \centering
    \includegraphics[width=0.45\textwidth]{slike/MaskP_curve.png}
    \hfill
    \includegraphics[width=0.45\textwidth]{slike/MaskR_curve.png}
    
    \vspace{0.5cm}
    \includegraphics[width=0.45\textwidth]{slike/MaskF1_curve.png}
    \hfill
    \includegraphics[width=0.45\textwidth]{slike/MaskPR_curve.png}
    \caption{Mask metrike: Preciznost (gore lijevo), Odziv (gore desno), F1 (dolje lijevo) i PR krivulja (dolje desno).}
\end{figure}

\section{IZDVAJANJE 3D ZNAČAJKI I PROCJENA POZE}

\subsection{Mapiranje i ekstrakcija oblaka točaka}
Nakon uspješne primjene obučenog 2D modela za segmentaciju instanci na RGB slikama iz složenih scena (npr. posuda s preklapajućim voćem), uslijedio je zadatak određivanja trodimenzionalne pozicije svake instance. Korištenjem dubinskih (Depth) podataka i intrinzičnih parametara kamere, izvršeno je RGB-D poravnanje. Svaki piksel iz 2D maske detektiranog objekta projiciran je u 3D prostor, čime je dobiven izolirani oblak točaka (engl. point cloud) za svaku pojedinu instancu voća u sceni.

\subsection{Određivanje centroida i usporedba metoda}
Za svaki ekstrahirani oblak točaka voća, izvršena je procjena 3D centroida (težišta objekta). Algoritam za izračun centroida implementiran je korištenjem dvije različite metode:
\begin{enumerate}
    \item \textbf{Pristup fitingom geometrijskog modela (Geometric model fitting):} Korištenje algoritama za prilagodbu aproksimativnih geometrijskih oblika (poput sfere ili elipsoida) na oblak točaka i pronalazak središta tako aproksimiranog tijela. Ovdje su u obzir uzeti samo podaci vidljivi sa strane senzora.
    \item \textbf{Ekstrakcija 3D granične kutije (3D Bounding Box):} Određivanje najmanjeg kvadra koji obuhvaća sve točke instance te izračun njegovog geometrijskog središta.
\end{enumerate}

Obje metode su uspoređene s obzirom na točnost pozicioniranja u uvjetima kada vidljivi oblak točaka predstavlja samo prednju plohu voća (zbog okluzije i kuta kamere). Rezultati su detaljno vizualizirani, a svaki student je analizu i evaluaciju odradio isključivo na klasi voća koja mu je dodijeljena.

\begin{figure}[h!]
    \centering
    \includegraphics[width=0.6\textwidth]{slike/3d_poza_pointcloud.png}
    \caption{Simulirana vizualizacija oblaka točaka s ekstrahiranom 3D graničnom kutijom, mapiranim centroidom i koordinatnim sustavom.}
    \label{fig:poza_pointcloud}
\end{figure}

\section{ZAKLJUČAK}
U ovom seminaru prikazan je cjelovit sustav za detekciju, segmentaciju i 3D lokalizaciju objekata u prostoru. Kroz zadatke prikupljanja prilagođenog skupa podataka, obučavanja modela i 3D obrade oblaka točaka demonstriran je kompletan inženjerski proces vizualne percepcije za robote. Rezultati primjene različitih metoda za procjenu težišta ukazuju na važnost odabira prikladnog algoritma s obzirom na nesavršenost i nepotpunost senzorskih očitavanja u stvarnom svijetu. Implementirani pipeline direktno doprinosi mogućnostima humanoidnih robota za složene zadatke manipulacije i samostalne interakcije s okolinom.

% --- LITERATURA ---
\newpage
\begin{thebibliography}{99}
\bibitem{fsb_materijali} Fakultet strojarstva i brodogradnje (2026). Materijali s predavanja i vježbi iz kolegija Humanoidna robotika.
\bibitem{roboflow} Roboflow (2026). Roboflow Documentation and Best Practices for Dataset Management.
\end{thebibliography}

\end{document}
